\documentclass[12pt,a4paper]{article}
\usepackage[utf8]{inputenc}
\usepackage[spanish]{babel}
\usepackage{amsmath,amssymb,amsfonts}
\usepackage{geometry}
\geometry{margin=2.5cm}

\title{\textbf{Compendio de Fórmulas de Acústica}}
\author{Física Acústica y Ondas}
\date{\today}

\begin{document}

\maketitle

\section*{Introducción}
La acústica es la rama de la física que estudia las ondas mecánicas que se propagan a través de la materia (gases, líquidos y sólidos). A continuación se presentan las fórmulas fundamentales de propagación, intensidad y efectos acústicos con sus unidades del SI.

\section{Ondas Sonoras y Propagación}

\subsection{Velocidad e Intensidad del Sonido}

\subsubsection*{1. Velocidad del Sonido e Intensidad Sonora [$c_s, I, L_I$]}
\paragraph*{Unidades básicas del SI:}
\begin{itemize}
    \item Velocidad del sonido [$c_s$]: $[\text{m}\cdot\text{s}^{-1}]$
    \item Intensidad sonora [$I$]: $[\text{kg}\cdot\text{s}^{-3}]$
    \item Nivel de intensidad sonora [$L_I$]: Adimensional $[1]$ (expresado en decibelios, $\text{dB}$)
\end{itemize}
\paragraph*{Explicación:} Define la rapidez con la que se propaga la onda de presión en un medio elástico y la potencia acústica transportada por unidad de área.
\paragraph*{Desarrollo y Variaciones:}
\begin{align*}
\text{Velocidad en fluidos (Ecuación de Laplace):} \quad & c_s = \sqrt{\frac{K}{\rho}} = \sqrt{\frac{\gamma R T}{M}} \\
\text{Velocidad en sólidos (barra delgada):} \quad & c_s = \sqrt{\frac{E}{\rho}} \\
\text{Intensidad sonora instantánea:} \quad & I = \langle p(t) v(t) \rangle = \frac{p_{\text{rms}}^2}{\rho c_s} \\
\text{Nivel de presión y de intensidad ($dB_{SPL}$):} \quad & L_I = 10 \log_{10} \left( \frac{I}{I_0} \right) = 20 \log_{10} \left( \frac{p_{\text{rms}}}{p_0} \right)
\end{align*}
donde $K$ es el módulo de compresibilidad, $E$ el módulo de Young, $\rho$ la densidad del medio, $p_{\text{rms}}$ la presión acústica eficaz y $I_0 = 10^{-12} \text{ W/m}^2$ la intensidad de referencia.
\paragraph*{Autor:} Pierre-Simon Laplace / Isaac Newton.

\section{Fenómenos Ondulatorios y Acústicos}

\subsection{Efecto Doppler}

\subsubsection*{2. Efecto Doppler Acústico [$f'$]}
\paragraph*{Unidades básicas del SI:} Frecuencia [$f'$]: $[\text{s}^{-1}]$ (Hertz)
\paragraph*{Explicación:} Aparente cambio de frecuencia de una onda producido por el movimiento relativo entre la fuente emisora y el observador.
\paragraph*{Desarrollo y Variaciones:}
\begin{align*}
\text{Fórmula general (fuente y observador en movimiento):} \quad & f' = f \left( \frac{c_s \pm v_o}{c_s \mp v_s} \right) \\
\text{Observador móvil, fuente en reposo ($v_s=0$):} \quad & f' = f \left( 1 \pm \frac{v_o}{c_s} \right) \\
\text{Fuente móvil, observador en reposo ($v_o=0$):} \quad & f' = f \left( \frac{c_s}{c_s \mp v_s} \right)
\end{align*}
donde $v_o$ es la velocidad del observador y $v_s$ la velocidad de la fuente respecto al medio.
\paragraph*{Autor:} Christian Doppler.

\end{document}
